% ============================================================================
% Fichier  : partie3/P03_C11_referentiels_controle.tex
% Titre    : Les Référentiels de Contrôle
% Auteur   : MINKA MI NGUIDJOÏ Thierry Emmanuel
% Version  : 1.0 -- novembre 2026
% Base   : ACSI C11, repris et développé pour INF3226
% ============================================================================
% Sources : COSO 2013 (évaluation : composantes et principes « présents et
% fonctionnels », déficience majeure) ; COBIT 2019 (40 objectifs, 5
% domaines ; copie de l'auteur, paraphrasée) et COBIT 4.1 (édition
% française, dossier Normes : 34 processus, 4 domaines) ; ISO/IEC
% 27001:2022 et 27002:2022 (93 mesures, 4 thèmes ; le dossier Normes
% contient l'édition 2013) ; NIST SP 800-53 rév. 5 (20 familles) et SP
% 800-53A rév. 5 (examiner, interroger, tester) ; guide d'audit des SI du
% CHAI (2014, neuf thèmes). Conformité à une norme = choix de fonction
% objectif (A0).
% ============================================================================

\chapter{Les Référentiels de Contrôle}
\label{chap:referentiels}

\epigraph{%
    « Le référentiel donne le critère.\\
    Il ne donne ni la preuve, ni la conclusion. »%
}{Adage des auditeurs}

\niveauChapitre{Fondamental}{%
    Maîtrise du Chapitre~\ref{chap:controle-interne} (COSO) et du
    Chapitre~\ref{chap:normes} (normes professionnelles)%
}

\begin{boxobjectifs}
\begin{itemize}
    \item Utiliser le COSO~2013 comme cadre d'évaluation d'un dispositif de
          contrôle interne.
    \item Utiliser COBIT~2019, et COBIT~4.1 encore cité dans de nombreux
          dossiers, comme catalogue d'objectifs de contrôle.
    \item Utiliser l'annexe~A de l'ISO/IEC~27001 et l'ISO/IEC~27002 comme
          catalogue de mesures de sécurité.
    \item Situer NIST SP~800-53 et le guide d'audit des SI du CHAI.
    \item Choisir les référentiels d'une mission, en sachant que la
          conformité à une norme est un choix d'objectif, pas un axiome.
\end{itemize}
\end{boxobjectifs}

\bigskip

% ============================================================================
\section*{Question directrice : par rapport à quoi juge-t-on~?}
% ============================================================================

Au Chapitre~\ref{chap:decomposition}, AUDIX CONSEIL a décomposé le domaine~D
de LiZbiZ en cinq sous-domaines et vingt-cinq processus. Chaque contrôle a
reçu une fiche d'assertion, avec un critère. Mais d'où vient ce critère~?
Faut-il exiger des mots de passe de huit caractères, comme la politique
interne, ou de douze, comme le propose le chef de mission~? Revoir les
comptes à privilèges chaque mois, ou chaque trimestre~?

\medskip
\textbf{Plan du chapitre.}
\begin{enumerate}
    \item Évaluer le contrôle interne avec le COSO (section~\ref{sec:c14-coso}).
    \item COBIT~2019, présenté pour lui-même, puis comme catalogue de critères (section~\ref{sec:c14-cobit}).
    \item L'ISO/IEC~27001 et l'ISO/IEC~27002 (section~\ref{sec:c14-iso}).
    \item Le NIST~: SP~800-53, SP~800-53A et le cadre CSF~2.0~; le guide du CHAI (section~\ref{sec:c14-autres}).
    \item Choisir les référentiels d'une mission, et ce que ce choix change aux mesures (section~\ref{sec:c14-choisir}).
\end{enumerate}

\separateur

% ============================================================================
\section{Évaluer le Contrôle Interne avec le COSO}
\label{sec:c14-coso}
% ============================================================================

\subsection{Présent, fonctionnel, intégré}

Le Chapitre~\ref{chap:controle-interne} a présenté les cinq composantes et
les dix-sept principes du COSO~2013. Le cadre sert aussi à \emph{évaluer}
un dispositif~\cite{coso2013ic}. Pour conclure qu'un contrôle interne est
efficace, trois conditions doivent être réunies~:
\begin{itemize}
    \item chaque composante et chaque principe pertinent est
          \textbf{présent}, c'est-à-dire conçu et mis en place~;
    \item chaque composante et chaque principe pertinent est
          \textbf{fonctionnel}, c'est-à-dire qu'il opère comme prévu~;
    \item les cinq composantes \textbf{agissent ensemble}, de manière
          intégrée.
\end{itemize}

\begin{boxdef}
\termecle{Déficience majeure}\index{déficience majeure} (COSO~2013)~:
déficience, ou combinaison de déficiences, qui réduit fortement la
probabilité que l'entité atteigne ses objectifs. Elle existe notamment
lorsqu'un principe pertinent n'est pas présent et fonctionnel. En présence
d'une déficience majeure, la direction ne peut pas conclure que son
contrôle interne est efficace.
\end{boxdef}

La règle est stricte~: dix-sept principes sur dix-sept. Elle ne se résume
pas à un taux de conformité, comme le montre l'évaluation de la BANQUE
SANAGA.

\begin{boxlecture}{Les 17 principes, SANAGA}
Pour l'évaluation exigée par la COBAC, Mme~EKAMBI cote chacun des dix-sept
principes « présent et fonctionnel »~: 12~\OUI{}, 3~\NON{}, 2~\NEK{}. Les
\NON{} portent sur le principe~9 (les changements significatifs, dont la
bascule de mars~2026, n'ont pas été analysés), le principe~11 (contrôles
généraux sur les changements) et le principe~14. Les \NEK{} portent sur les
principes~16 et~17~: aucune pièce n'établit que les évaluations ont eu lieu.

\medskip
\textbf{Constatation.} $\mu \approx 0{,}71$~; $U \approx 0{,}12$~;
$H \approx 1{,}16$~bit. Selon le quadrant du cours, le domaine est en
conformité claire.

\medskip
\textbf{Analyse} (confiance élevée). Selon le COSO, le dispositif n'est
pas efficace~: trois principes ne sont pas présents et fonctionnels, ce
sont trois déficiences majeures, et deux autres principes ne peuvent pas
être établis. Le référentiel impose sa propre règle de décision, plus
exigeante que le seuil de $\mu = 0{,}5$. L'auditeur applique les deux et
les rapporte toutes deux~: les mesures disent \emph{où} l'on en est, la
règle du COSO dit \emph{si} l'on peut conclure.

\medskip
\hfill\lecture{0{,}71}{0{,}12}{1{,}16}
\end{boxlecture}

\begin{boxmetier}{Contrôleur interne}
\textbf{Préparer l'évaluation COSO d'un dispositif}
\begin{enumerate}
    \item Pour chaque principe, lister les contrôles qui le mettent en
          œuvre, et le responsable de chacun.
    \item Pour chaque principe, rassembler la preuve qu'il est présent
          (procédure, organigramme, charte) et la preuve qu'il fonctionne
          (traces de l'année).
    \item Coter chaque principe \OUI{}, \NON{} ou \NEK{}~; un principe sans
          preuve de fonctionnement est \NEK{}, pas \OUI{}.
    \item Qualifier chaque déficience et dire si elle est majeure.
    \item Conclure selon la règle du COSO, puis rapporter $\mu$, $U$ et $H$
          pour situer l'écart.
\end{enumerate}
\end{boxmetier}

% ============================================================================
\section{COBIT : un Catalogue d'Objectifs de Contrôle}
\label{sec:c14-cobit}
% ============================================================================

\subsection{COBIT 2019}

Le cours GCI (MSN-2S3-MSI03) présente COBIT~2019 comme référentiel de
gouvernance. L'auditeur l'utilise autrement~: comme un catalogue de
pratiques, où il puise des critères. COBIT~2019 compte quarante objectifs
de gouvernance et de gestion, répartis en cinq
domaines~\cite{isaca2018cobitintro,isaca2018cobitobjectifs}.

\begin{tableaucours}{cL{4.4cm}cL{5.4cm}}{%
    \tableauHeader{Code} &
    \tableauHeader{Domaine} &
    \tableauHeader{Obj.} &
    \tableauHeader{Exemples utiles à l'auditeur}
}
EDM & Évaluer, diriger et surveiller & 5 & EDM03, optimisation des risques \\
APO & Aligner, planifier et organiser & 14 & APO12, risques~; APO13,
sécurité \\
BAI & Bâtir, acquérir et implémenter & 11 & BAI06, changements \\
DSS & Délivrer, servir et soutenir & 6 & DSS04, continuité~; DSS05,
services de sécurité \\
MEA & Surveiller, évaluer et apprécier & 4 & MEA02, contrôle interne~;
MEA03, conformité~; MEA04, assurance \\
\end{tableaucours}
\vspace{-0.8em}
\begin{center}{\small\itshape Tableau~14.1. Les cinq domaines de COBIT~2019
(intitulés traduits)}\end{center}

Chaque objectif se décline en pratiques de gestion, elles-mêmes détaillées
en activités. La pratique est le bon niveau pour un critère d'audit~: par
exemple, la pratique DSS05.04 traite de la gestion des identités et des
accès logiques, et fournit les critères des sous-domaines D.2 et D.4 de
LiZbiZ.

\begin{boxinfo}
\textbf{Maturité ou capacité~?} COBIT~2019 évalue la \emph{capacité} des
processus sur une échelle de 0 à~5, activité par activité. Le
Chapitre~\ref{chap:conclure} a rappelé que le niveau de maturité du cours,
tiré de $\mu$, est une convention distincte~: les deux échelles ne se
convertissent pas l'une dans l'autre.
\end{boxinfo}

\subsection{COBIT 2019 pour lui-même}

Un élève ingénieur qui n'a pas suivi le cours de gouvernance doit pouvoir
lire COBIT seul. Le référentiel repose sur trois idées~\cite{isaca2018cobitintro}.
\begin{enumerate}
    \item \textbf{Gouvernance et gestion sont distinctes.} La gouvernance
          (domaine EDM) évalue, oriente et surveille~; elle relève du
          conseil. La gestion (APO, BAI, DSS, MEA) planifie, construit,
          exploite et surveille~; elle relève de la direction.
    \item \textbf{Un système de gouvernance a plusieurs composants}~:
          processus, structures organisationnelles, principes et
          politiques, information, culture et comportements, compétences,
          services et infrastructures. Auditer un objectif COBIT, ce n'est
          pas seulement vérifier un processus~: c'est aussi vérifier qu'une
          structure en est responsable et qu'une politique le fonde.
    \item \textbf{Le système se conçoit selon le contexte}~: des
          \emph{facteurs de conception} (stratégie, profil de risque,
          taille, rôle de l'informatique, réglementation\ldots) font varier
          l'importance de chaque objectif d'une organisation à l'autre.
\end{enumerate}

\begin{boxinfo}
\textbf{Ce que l'auditeur en retient.} Pour la BANQUE SANAGA, un facteur
de conception domine~: l'exigence réglementaire (COBAC). Il rend
prioritaires MEA03 (conformité), BAI06 (changements) et DSS04
(continuité)~; ce sont les objectifs où l'auditeur ira chercher ses
critères en premier.
\end{boxinfo}

\subsection{COBIT 4.1, encore très cité}

De nombreux dossiers d'audit, et l'ancien cours de l'auteur, se réfèrent à
COBIT~4.1, publié en 2007 et disponible en français~\cite{itgi2007cobit41}.
Il comptait trente-quatre processus en quatre domaines (PO, AI, DS, ME),
chacun décliné en objectifs de contrôle, et proposait des modèles de
maturité de 0 à~5. Il faut savoir passer de l'un à l'autre.

\begin{tableaucours}{L{4.6cm}L{3.6cm}L{4.0cm}}{%
    \tableauHeader{Sujet} &
    \tableauHeader{COBIT 4.1} &
    \tableauHeader{COBIT 2019}
}
Sécurité des systèmes & DS5 & APO13, DSS05 \\
Gestion des changements & AI6 & BAI06 \\
Continuité du service & DS4 & DSS04 \\
Surveillance du contrôle interne & ME2 & MEA02 \\
Conformité aux exigences externes & ME3 & MEA03 \\
\end{tableaucours}
\vspace{-0.8em}
\begin{center}{\small\itshape Tableau~14.2. Quelques correspondances entre
COBIT~4.1 et COBIT~2019}\end{center}

% ============================================================================
\section{ISO/IEC 27001 et 27002 : les Mesures de Sécurité}
\label{sec:c14-iso}
% ============================================================================

L'ISO/IEC~27001 fixe les exigences d'un SMSI~; son annexe~A énumère des
mesures de sécurité, que l'ISO/IEC~27002 détaille et
explique~\cite{iso27001_2022,iso27002_2022}. Depuis l'édition de 2022,
l'annexe compte 93~mesures en quatre thèmes~: organisationnelles~(37),
humaines~(8), physiques~(14) et technologiques~(34). L'édition de 2013, que
contient le dossier \texttt{Normes}, en comptait 114 en quatorze domaines~:
un critère tiré de l'ancienne numérotation doit être converti.

Chaque mesure de l'édition 2022 porte des \emph{attributs}~: type de mesure
(préventive, détective, corrective), propriétés de sécurité, concepts de
cybersécurité, capacités opérationnelles, domaines de sécurité. Ils
permettent à l'auditeur de filtrer le catalogue~: par exemple, ne retenir
que les mesures détectives d'un domaine.

\begin{boiteRemarque}{L'annexe A est-elle une grille d'audit~?}
Non. Recopier les 93~mesures en autant de questions donne une grille qui
viole la borne cognitive (Chapitre~\ref{chap:decomposition}) et pose des
questions sans objet. L'entité choisit ses mesures dans sa \emph{déclaration
d'applicabilité}, en justifiant chaque exclusion. L'auditeur part de sa
décomposition, puis va chercher dans l'annexe le critère de chaque contrôle,
et vérifie au passage que les exclusions déclarées sont justifiées.
\end{boiteRemarque}

% ============================================================================
\section{NIST SP 800-53 et le Guide du CHAI}
\label{sec:c14-autres}
% ============================================================================

\subsection{NIST SP 800-53}

La publication spéciale SP~800-53 du NIST (révision~5, 2020) est un
catalogue de mesures de sécurité et de protection de la vie privée,
réparties en vingt familles (contrôle d'accès AC, journalisation AU,
gestion de configuration CM, continuité CP, identification et
authentification IA, sécurité du personnel PS, etc.)~\cite{nist2020sp80053}.
Elle s'impose aux systèmes fédéraux américains, et sert ailleurs de
catalogue de référence, souvent plus précis que l'ISO/IEC~27002.

Son intérêt pour l'auditeur tient surtout à la publication compagne
SP~800-53A, qui donne pour chaque mesure des procédures d'évaluation selon
trois méthodes~: \emph{examiner}, \emph{interroger}, \emph{tester}~\cite{nist2022sp80053a}.
Elles correspondent aux tests DOC, ENT et TEST du
Chapitre~\ref{chap:preuve}.

\subsection{Le cadre de cybersécurité du NIST (CSF 2.0)}

Le NIST publie aussi, depuis 2014, un cadre de cybersécurité plus
synthétique, révisé en février~2024 (version~2.0)~\cite{nist2024csf}. Il
organise les résultats attendus en six \emph{fonctions}.

\begin{tableaucours}{L{2.6cm}L{9.6cm}}{%
    \tableauHeader{Fonction} &
    \tableauHeader{Résultats attendus (reformulés)}
}
Gouverner & Stratégie, rôles, politique et supervision des risques de
cybersécurité (fonction ajoutée en 2024) \\
Identifier & Connaître les actifs, les risques, les fournisseurs \\
Protéger & Mettre en place les protections (identités, données,
plateformes) \\
Détecter & Repérer les événements et les anomalies \\
Répondre & Traiter un incident détecté \\
Rétablir & Restaurer les actifs et les opérations \\
\end{tableaucours}
\vspace{-0.8em}
\begin{center}{\small\itshape Tableau~14.3. Les six fonctions du NIST
CSF~2.0}\end{center}

Le CSF ne fixe pas de critères détaillés~: il renvoie à d'autres
catalogues (dont la SP~800-53 et l'ISO/IEC~27002). Il sert à l'auditeur à
vérifier la \emph{couverture} d'un dispositif~: une organisation qui
protège bien mais ne sait ni détecter ni rétablir présente un défaut que
la seule lecture de l'annexe~A ne révèle pas toujours.

\subsection{Le guide d'audit des SI du CHAI}

Le guide du Comité d'harmonisation de l'audit interne (CHAI), élaboré en
2014 par un groupe de travail interministériel français, est un guide
pratique, en français, organisé par thèmes d'audit~\cite{chai2014guide}~:
missions généralistes, pilotage des SI, sécurité, production, applications
en service, support et parc, fonction « études », projets, marchés
informatiques. Il ne prétend pas être le seul référentiel~: il propose,
pour chaque thème, des points de contrôle. La Partie~\ref{part:pratique}
suit largement ce découpage.

% ============================================================================
\section{Choisir les Référentiels d'une Mission}
\label{sec:c14-choisir}
% ============================================================================

\subsection{Un même contrôle, quatre critères}

Les référentiels se recouvrent. Pour un même contrôle, l'auditeur trouve
plusieurs critères possibles~; il en retient un, et le cite dans la fiche
d'assertion.

\begin{tableaucours}{L{4.0cm}L{1.8cm}L{1.8cm}L{1.8cm}L{2.2cm}}{%
    \tableauHeader{Contrôle} &
    \tableauHeader{COSO} &
    \tableauHeader{COBIT 2019} &
    \tableauHeader{ISO/IEC 27002} &
    \tableauHeader{NIST 800-53}
}
D.2.1 Désactivation des comptes des partants & Principe~11 & DSS05.04 &
5.18 & AC-2, PS-4 \\
D.4.2 Revue des comptes à privilèges & Principe~11 & DSS05.04 & 8.2 &
AC-2, AC-6 \\
B Sauvegardes & Principe~11 & DSS04 & 8.13 & CP-9 \\
Gestion des changements (SANAGA) & Principes~9 et~11 & BAI06 & 8.32 &
CM-3 \\
\end{tableaucours}
\vspace{-0.8em}
\begin{center}{\small\itshape Tableau~14.4. Correspondances entre
référentiels pour quatre contrôles (à vérifier dans l'édition en vigueur)}\end{center}

\subsection{Les critères de choix}

\begin{boxprocedure}
\textbf{Choisir les référentiels.} Dans l'ordre~:
\begin{enumerate}
    \item \textbf{Ce qui s'impose}~: loi, règlement, exigence du régulateur
          (COBAC pour la BANQUE SANAGA, régime financier de l'État pour
          l'ANPP, Chapitre~\ref{chap:droit}).
    \item \textbf{Ce que l'entité s'est engagée à respecter}~: ses propres
          politiques, un contrat, une certification revendiquée.
    \item \textbf{Ce que le commanditaire demande}~: une comparaison avec un
          référentiel international, pour se situer.
    \item \textbf{Ce que l'auditeur juge nécessaire}, en le justifiant~:
          un critère plus exigeant que la politique interne, lorsque le
          risque le demande.
\end{enumerate}
Le choix est écrit dans la lettre de mission ou dans le plan d'audit,
discuté avec le commanditaire, et repris dans le rapport.
\end{boxprocedure}

\subsection{Le référentiel, une fonction objectif}

L'article A0 tire de ce choix une conséquence théorique~: la conformité à
un référentiel n'est pas un fondement de l'audit, c'est le choix d'une
\emph{fonction objectif}~\cite{minka2026auditologie}. Changer de référentiel
change ce que l'on appelle « conforme »~; cela ne change pas ce que l'on
sait. La lecture suivante le montre.

\begin{boxlecture}{Domaine D.3, deux critères}
Au Chapitre~\ref{chap:decomposition}, la question du sous-domaine D.3
(authentification), posée aux quatre interlocuteurs, est cotée d'après la
politique interne de LiZbiZ~: 4~\OUI{} sur 4~réponses. Le chef de mission
propose de retenir des critères plus exigeants, tirés de la mesure~8.5 de
l'ISO/IEC~27002~: mots de passe de douze caractères, et double
authentification pour \emph{tout} accès distant, et non pour les seuls
administrateurs. Les quatre réponses sont cotées de nouveau sur les mêmes
preuves.

\medskip
\begin{tableaucours}{L{4.6cm}cccccc}{%
    \tableauHeader{Critère} &
    \tableauHeader{\OUI{}} & \tableauHeader{\NON{}} & \tableauHeader{\NEK{}} &
    \tableauHeader{$\mu$} & \tableauHeader{$U$} & \tableauHeader{$H$}
}
Politique interne de LiZbiZ & 4 & 0 & 0 & 1,00 & 0,00 & 0,00 \\
Critères tirés de l'ISO/IEC~27002 & 2 & 2 & 0 & 0,50 & 0,00 & 1,00 \\
\end{tableaucours}

\medskip
\textbf{Constatation.} Le taux de conformité passe de 1,00 à 0,50~; le
taux d'incertitude reste nul.

\medskip
\textbf{Analyse} (confiance élevée). Le référentiel a déplacé $\mu$, pas
$U$~: l'état de connaissance de l'auditeur n'a pas changé, seule la cible a
changé. C'est pourquoi le rapport doit nommer le critère de chaque
conclusion~; et c'est pourquoi l'auditeur ne change pas de référentiel en
cours de mission sans l'accord écrit du commanditaire.

\medskip
\hfill\lecture{0{,}50}{0{,}00}{1{,}00}
\end{boxlecture}

\begin{boxalert}
\textbf{Violation typique.} Annoncer dans le rapport « conformité à
l'ISO/IEC~27001 à 62~\% » à partir d'une grille construite sur la politique
interne de l'entité.

\smallskip
\textbf{Règle~:} un taux de conformité n'a de sens que rapporté au critère
qui l'a produit. On écrit « 62~\% des contrôles conformes à la politique
interne, évalués selon les critères du plan d'audit », et l'on ne parle de
conformité à une norme que si les critères en sont tirés.
\end{boxalert}

\begin{boxmetier}{Gestionnaire de risques}
\textbf{Tenir une table de correspondance}
\begin{enumerate}
    \item Recenser les référentiels qui s'imposent à l'entité, et ceux
          qu'elle a choisis.
    \item Définir des \emph{contrôles communs}, chacun relié aux exigences
          de plusieurs référentiels (COSO, COBIT, ISO/IEC~27002, COBAC).
    \item Tester chaque contrôle commun une seule fois, et réutiliser la
          preuve pour chaque référentiel.
    \item Mettre à jour la table à chaque nouvelle édition d'un
          référentiel, en signalant les numéros qui changent.
    \item Remettre la table à l'auditeur dès la préparation de la mission.
\end{enumerate}
\end{boxmetier}

% ============================================================================

\separateur

\begin{boxresume}
\textbf{Ce qu'il faut retenir de ce chapitre~:}
\begin{itemize}
    \item \textbf{Le COSO} évalue un dispositif~: chaque principe présent et
          fonctionnel, les composantes intégrées~; un seul principe
          défaillant est une déficience majeure.
    \item \textbf{COBIT~2019} (40 objectifs, 5 domaines) est un catalogue de
          pratiques~; COBIT~4.1 (34 processus) reste très cité, et se
          convertit.
    \item \textbf{L'ISO/IEC~27002:2022} détaille 93~mesures en quatre
          thèmes~; l'annexe~A n'est pas une grille d'audit.
    \item \textbf{NIST SP~800-53} est un catalogue détaillé~; la SP~800-53A
          donne des procédures d'évaluation (examiner, interroger, tester).
    \item \textbf{Le guide du CHAI} propose des points de contrôle par thème
          d'audit.
    \item \textbf{Le choix du référentiel} est celui d'une fonction
          objectif~: il déplace $\mu$, pas $U$, et doit être écrit, discuté
          et cité.
    \item \textbf{COBIT~2019} distingue gouvernance et gestion et s'adapte par facteurs de conception~; le \textbf{NIST CSF~2.0} vérifie la couverture en six fonctions, dont « Gouverner ».
\end{itemize}
\end{boxresume}

\bigskip

\begin{boxexo}[title={\ding{45}\quad Exercice~14.1. Trouver le critère \niveaubase}]
Pour chaque contrôle, indiquez l'objectif COBIT~2019 et la mesure de
l'ISO/IEC~27002:2022 où chercher le critère.
\begin{enumerate}[(a)]
    \item Les sauvegardes du centre de Douala sont testées chaque trimestre.
    \item Toute modification de l'application intégrée est approuvée avant
          sa mise en production.
    \item Les comptes des agents partis sont désactivés sous 48~heures.
    \item Les droits d'administration sont revus chaque mois.
\end{enumerate}
\end{boxexo}

\begin{boxexo}[title={\ding{45}\quad Exercice~14.2. Conclure selon le COSO \niveaumoyen}]
Reprenez la lecture « Les 17 principes, SANAGA ».
\begin{enumerate}[(a)]
    \item Pourquoi un taux de conformité de 0,71 ne permet-il pas de
          conclure à l'efficacité du dispositif~?
    \item Quelles pièces lèveraient les deux \NEK{} des principes~16 et~17~?
    \item Si ces deux \NEK{} étaient levés en \OUI{}, la conclusion
          changerait-elle~? Justifiez.
    \item Rédigez la conclusion de l'évaluation destinée à la COBAC.
\end{enumerate}
\end{boxexo}

\begin{boxexo}[title={\ding{45}\quad Exercice~14.3. Choisir les référentiels de l'ANPP \niveauexpert}]
La Chambre des comptes prépare l'audit des paiements de l'ANPP, dont les
comptes d'agents partis encore actifs.
\begin{enumerate}[(a)]
    \item Quels référentiels s'imposent, et lesquels sont choisis~?
          Classez-les selon la procédure du chapitre.
    \item Pour le contrôle « les comptes des agents partis sont désactivés
          sous 48~heures », l'ANPP n'a aucune procédure écrite. Quel critère
          retenir, et comment le justifier~?
    \item Montrez, à l'aide de la lecture du domaine D.3, pourquoi le rapport
          doit nommer le critère de chaque conclusion.
    \item En quoi le choix d'un référentiel est-il, au sens de l'article A0,
          un choix de fonction objectif~?
\end{enumerate}
\end{boxexo}

\bigskip

\begin{boxexo}[title={\ding{45}\quad Exercice~14.4. Gouvernance ou gestion~? \niveaubase}]
Rangez dans le bon domaine de COBIT~2019~: (a) le conseil approuve
l'appétence au risque informatique~; (b) la DSI teste un changement avant
sa mise en production~; (c) l'audit interne évalue le contrôle interne~;
(d) l'exploitation restaure une sauvegarde~; (e) la direction établit le
schéma directeur.
\end{boxexo}

\begin{boxexo}[title={\ding{45}\quad Exercice~14.5. Les fonctions du CSF \niveaubase}]
Classez les constats suivants de LiZbiZ dans les fonctions du NIST
CSF~2.0~: comptes d'agents partis actifs~; aucun inventaire des
applications~; journaux jamais revus~; restaurations jamais testées~;
aucune politique de sécurité approuvée.
\end{boxexo}

\begin{boxexo}[title={\ding{45}\quad Exercice~14.6. Convertir une ancienne grille \niveaumoyen}]
Un dossier d'audit de 2019 de la BANQUE SANAGA cite « COBIT~4.1 DS5 » et
« ISO/IEC~27002:2013, 9.2 ». Donnez les références à utiliser en 2026 et
expliquez pourquoi le rapport ne peut pas reprendre telles quelles les
anciennes références.
\end{boxexo}

\begin{boxexo}[title={\ding{45}\quad Exercice~14.7. La déclaration d'applicabilité \niveaumoyen}]
La déclaration d'applicabilité de DIGICORE AFRIQUE exclut la mesure de
l'ISO/IEC~27002 sur la sécurité des développements, au motif que
« DIGICORE ne développe pas ». Or DIGICORE paramètre et modifie le système
central de la banque. Quelle est votre analyse, et quelle cote donnez-vous
au contrôle « les exclusions sont justifiées »~?
\end{boxexo}

\begin{boxexo}[title={\ding{45}\quad Exercice~14.8. Examiner, interroger, tester \niveaumoyen}]
Pour la mesure « revue des droits d'accès », rédigez une procédure
d'évaluation dans chacune des trois méthodes de la SP~800-53A, et
associez-la au type de test du Chapitre~\ref{chap:preuve}.
\end{boxexo}

\begin{boxexo}[title={\ding{45}\quad Exercice~14.9. Le référentiel déplace µ \niveauexpert}]
Reprenez le sous-domaine D.3 de LiZbiZ (authentification), coté 4/0/0
avec la politique interne et 2/2/0 avec un critère plus exigeant.
Expliquez pourquoi le changement de référentiel déplace $\mu$ mais pas
$U$, et dites quel critère vous retenez pour le rapport, en justifiant par
les critères de choix du chapitre.
\end{boxexo}

\begin{boxexo}[title={\ding{45}\quad Exercice~14.10. Un référentiel pour LiZbiZ \niveauexpert}]
Proposez, pour les cinq domaines de LiZbiZ, le référentiel principal et
un référentiel complémentaire, et justifiez. Indiquez pour quel domaine le
COSO est indispensable, et pour lequel le NIST CSF apporte le plus.
\end{boxexo}

\bigskip

\begin{boxinfo}
\textbf{Points de contrôle de l'élève.} Je sais~: conclure sur un
principe COSO (présent, fonctionnel, intégré)~; citer les cinq domaines de
COBIT~2019 et leur nombre d'objectifs~; passer de COBIT~4.1 à COBIT~2019~;
dire combien de mesures compte l'annexe A de 2022~; nommer les six
fonctions du CSF~2.0~; choisir et justifier un référentiel.
\end{boxinfo}

\begin{boxinfo}
\textbf{Pour aller plus loin.}
\begin{itemize}
    \item COSO, \textit{Internal Control -- Integrated
          Framework}~\cite{coso2013ic}.
    \item ISACA, \textit{COBIT 2019 Framework}~\cite{isaca2018cobitintro,isaca2018cobitobjectifs}.
    \item ISO/IEC~27002:2022~\cite{iso27002_2022}.
    \item NIST, SP~800-53 rév.~5 et SP~800-53A~\cite{nist2020sp80053,nist2022sp80053a}~;
          \textit{Cybersecurity Framework 2.0}~\cite{nist2024csf}.
\end{itemize}
\end{boxinfo}

\begin{boxinfo}
\textbf{Transition.} Les référentiels se choisissent~; la loi s'impose. Le
Chapitre~\ref{chap:droit} présente le cadre juridique camerounais et
africain de l'audit et du contrôle des SI~: loi n°~2010/012, règlement
COBAC, AUSCGIE et contrôle des finances publiques.
\end{boxinfo}
